% !TeX program = lualatex
% !TeX TXS-program:compile = txs:///lualatex/[--shell-escape]
% The reference picture's notes, with what the paper overhaul added: a crumpled sheet,
% a coffee stain, underlined headings, a margin note and hand-drawn arrows.
% Build:  bin\cacography.cmd qm-handled
\documentclass{article}
\usepackage[
  paper=notebook,
  ink=blue,
  sloppiness=0.5,
  seed=123,
  underline=headings,
  background=walnut,
  light=window,
  framing=tight,
  props=none,
]{cacography}
\usepackage{amsmath}

\title{Quantum Mechanics (very messy notes)}
\author{A. Student}
\date{}

\begin{document}
\cacographypage{crumple=0.8, stains=coffee, edges=worn, backside=writing}
\maketitle

\section{Schr\"odinger Equation}
The time-dependent Schr\"odinger equation\cacographynote*{maybe add time-indep. case here?} is
\begin{equation}
  i\hbar \frac{\partial}{\partial t}\psi(\mathbf{r},t)
  = -\frac{\hbar^2}{2m}\nabla^2\psi(\mathbf{r},t) + V(\mathbf{r},t)\psi(\mathbf{r},t)
\end{equation}
For a stationary state we can write $\psi(\mathbf{r},t) = \phi(\mathbf{r})e^{-iEt/\hbar}$,
which gives
\begin{equation}
  -\frac{\hbar^2}{2m}\nabla^2\phi(\mathbf{r}) + V(\mathbf{r})\phi(\mathbf{r}) = E\phi(\mathbf{r}).
\end{equation}

\section{Some remarks}
Note:
\begin{itemize}
  \item $V(\mathbf{r})$ can be complicated.
  \item Boundary conditions matter!
  \item Check normalization: $\int |\psi|^2 \, d^3r = 1$.
  \item Stationary $\Rightarrow$ $|\psi|^2$ does not change with time.
\end{itemize}

\end{document}
